\chapter{Conclusão}
\label{chap:conclusao}

Este trabalho teve como objetivo desenvolver e avaliar um sistema de reconhecimento facial capaz de atuar como alternativa aos cartões de identificação utilizados em ambientes institucionais. Dando continuidade aos estudos realizados no \acs{PFC} I, foi implementado um sistema completo de reconhecimento facial em tempo real, contemplando captura por câmera, detecção facial, extração de \textit{embeddings}, armazenamento seguro das informações biométricas, identificação de usuários cadastrados e integração com um sistema institucional para validação automática de credenciais de acesso.

Os experimentos realizados ao longo da pesquisa permitiram analisar o comportamento dos \textit{embeddings} faciais gerados pelo modelo adotado e avaliar o desempenho do sistema em diferentes condições de operação. Os resultados demonstraram que os \textit{embeddings} organizam o espaço vetorial de forma consistente, mantendo proximidade entre amostras pertencentes à mesma identidade e separação adequada entre indivíduos distintos. A análise das métricas biométricas, em especial \acs{FAR} (\textit{False Acceptance Rate}) e \acs{FRR} (\textit{False Rejection Rate}), permitiu determinar valores apropriados de \textit{threshold}, estabelecendo um equilíbrio entre segurança e usabilidade compatível com os requisitos de sistemas de autenticação biométrica.

Além disso, os experimentos de identificação no cenário 1:N demonstraram que o sistema é capaz de recuperar corretamente a identidade mais provável a partir da base de usuários cadastrados.

A solução desenvolvida utilizou o detector facial \acs{MTCNN} para localização e alinhamento das faces e o modelo \textit{FaceNet} (\textit{InceptionResnetV1}) para geração dos \textit{embeddings} faciais. Como forma de aumentar a confiabilidade do reconhecimento, foram incorporados mecanismos de avaliação da qualidade das imagens capturadas, cadastro multipose e votação temporal baseada em múltiplos quadros consecutivos. Adicionalmente, os \textit{embeddings} foram armazenados localmente em banco de dados SQLite com criptografia, contribuindo para a proteção das informações biométricas dos usuários.

Além da validação experimental, o sistema foi integrado e testado em um ambiente institucional real para verificação de credenciais de acesso. Durante os testes realizados, a solução demonstrou funcionamento satisfatório, permitindo a identificação dos usuários e a recuperação automática das informações associadas às suas credenciais. Os resultados obtidos reforçam a viabilidade prática da utilização do reconhecimento facial como mecanismo de autenticação em aplicações acadêmicas e institucionais.

Dessa forma, os objetivos propostos para o trabalho foram alcançados. Foi possível validar experimentalmente a abordagem baseada em \textit{embeddings}, implementar um sistema funcional de reconhecimento facial, armazenar e gerenciar identidades de forma segura, realizar identificação em tempo real e integrar o reconhecimento biométrico a um fluxo de autenticação utilizado em ambiente institucional. Assim, a pesquisa avançou da validação conceitual apresentada no \acs{PFC} I para uma implementação prática e operacional da proposta.

Apesar dos resultados positivos, algumas limitações permanecem. O sistema ainda depende de condições adequadas de iluminação e posicionamento do usuário, não possui mecanismos avançados de detecção de vivacidade (\textit{liveness detection}) e pode demandar ajustes de parâmetros para diferentes ambientes e equipamentos. Além disso, a estimativa de pose facial empregada utiliza heurísticas geométricas simplificadas, podendo ser aprimorada por abordagens mais sofisticadas.

Como trabalhos futuros, destacam-se a implementação de mecanismos de detecção de vivacidade para aumentar a segurança contra tentativas de fraude, a criação de ferramentas de auditoria e monitoramento dos eventos de reconhecimento, a modularização da arquitetura do sistema, a realização de avaliações em bases de usuários maiores e mais diversificadas e a integração com dispositivos físicos de controle de acesso. Também se destaca a expansão da solução para novos cenários de aplicação, incluindo a identificação automática de veículos por meio de visão computacional, permitindo a integração entre o reconhecimento de pessoas e o controle de acesso veicular em ambientes institucionais.

Em síntese, os resultados obtidos demonstram que o reconhecimento facial baseado em \textit{embeddings} constitui uma alternativa tecnicamente viável para autenticação e controle de acesso. A combinação entre os experimentos realizados e a implementação do protótipo funcional permitiu validar tanto a eficácia da abordagem quanto sua aplicabilidade prática, evidenciando o potencial da solução para aplicações reais e fornecendo uma base sólida para futuras evoluções da plataforma.
